% !TeX root = ../../Book5.tex
% 纲要标题：## 附录 B　代数叠、导出几何与几何化
% 纲要位置：几何学计划纲要.md，第 6352 行。
% 正文状态：未撰写；本文件目前是框架。

% 纲要内容开始
% 
% > 在正文的具体构造之后延伸，不再次仅重复名词定义。
% 
% * **商叠与层范畴**
%   * BG 和 [X/G] 的显式对象
%   * 拟凝聚层与等变模
%   * 稳定子和粗模的失真
% 
% * **导出局部模型**
%   * 自交的高阶 Tor
%   * 余切复形和切复形的对偶条件
%   * 完美障碍理论的动机
%   * 一般导出叠的新语言
% 
% * **几何表示的低维模型**
%   * 旗簇上的线丛
%   * 联络、单值和解层
%   * 推拉函子与卷积
%   * 一般等价定理的范围
% 
% ---
% 
% 纲要内容结束

\chapter{代数叠、导出几何与几何化}
\label{g5:app:B}
\BookChapterNumber{g5:app:B}

待撰写

\section{商叠与层范畴}
\label{g5:sec:B01}
待撰写

\subsection{BG 和 [X/G] 的显式对象}
\label{g5:subsec:B01:01}
待撰写

\subsection{拟凝聚层与等变模}
\label{g5:subsec:B01:02}
待撰写

\subsection{稳定子和粗模的失真}
\label{g5:subsec:B01:03}
待撰写

\section{导出局部模型}
\label{g5:sec:B02}
待撰写

\subsection{自交的高阶 Tor}
\label{g5:subsec:B02:01}
待撰写

\subsection{余切复形和切复形的对偶条件}
\label{g5:subsec:B02:02}
待撰写

\subsection{完美障碍理论的动机}
\label{g5:subsec:B02:03}
待撰写

\subsection{一般导出叠的新语言}
\label{g5:subsec:B02:04}
待撰写

\section{几何表示的低维模型}
\label{g5:sec:B03}
待撰写

\subsection{旗簇上的线丛}
\label{g5:subsec:B03:01}
待撰写

\subsection{联络、单值和解层}
\label{g5:subsec:B03:02}
待撰写

\subsection{推拉函子与卷积}
\label{g5:subsec:B03:03}
待撰写

\subsection{一般等价定理的范围}
\label{g5:subsec:B03:04}
待撰写
